\documentclass[11pt,a4paper]{article}
\usepackage[margin=2.4cm]{geometry}
\usepackage[T1]{fontenc}
\usepackage{amsmath,amssymb,booktabs,graphicx,hyperref}
\newcommand{\Call}{\operatorname{Call}}
\newcommand{\snap}{../results/BTC_20261114T223158Z}

\title{Replicating SANOS on Deribit BTC options}
\author{Gabriel Borrego Rold\'an}
\date{November 2026}

\begin{document}
\maketitle

\begin{abstract}
We reimplement the production model of SANOS (Buehler, Horvath, Kratsios, Limmer and Saqur,
arXiv:2601.11209v4) and test seven of its claims, first on synthetic quotes from a known
model and then on two snapshots of the Deribit BTC option book. Five claims hold, one holds
on a smaller problem than the paper's, and one, sensible extrapolation beyond the quoted
strikes, is not reproduced. The real-data evidence is thin: the two snapshots are four
minutes apart, so they show one market state.
\end{abstract}

\section{The model}

Write prices in the ``pure'' coordinates of the paper: strikes divided by the forward, so
that a call price is $C(T,k)=\mathbb{E}[(Z_T-k)^+]$ for a positive martingale $Z$ with
$Z_0=1$. For expiries $T_1<\dots<T_M$, SANOS sets
\begin{equation}
  \hat C(T_j,k)=\sum_{i=1}^{N_j} q_j^i\,\Call\bigl(K_j^i,k,\eta V_j\bigr),
\end{equation}
where $\Call(s,k,v)$ is the Black--Scholes price with spot $s$ and total variance $v$,
the $K_j^i$ are model strikes, $V_j$ are increasing backbone variances, $\eta\in[0,1)$
is a smoothness parameter, and each $q_j$ is a probability vector with
$\sum_i q_j^iK_j^i=1$. The price is that of $Z_j=X_jY_j$ with $X_j\sim q_j$ discrete and
$Y_j$ an independent unit-mean lognormal: a discrete distribution blurred multiplicatively.
For $\eta=0$ it is linear interpolation of call prices.

The surface is free of calendar arbitrage if the discrete call prices of the $q_j$ increase
in $j$, which is linear in $q$:
$\sum_i q_j^i(K_j^i-K)^+\ge\sum_i q_{j-1}^i(K_{j-1}^i-K)^+$ at the strikes $K=K_j^\ell$.
The paper calls this $\omega=0$; imposing the inequality on the smoothed prices instead is
$\omega=1$ and guarantees nothing between strikes. Since model prices are linear in $q$,
fitting bid/ask quotes is a linear program. We use the paper's recommended objective: with
weights $1/(\text{ask}-\text{bid})$, the distance outside bid/ask plus $10^{-8}$ times the
distance to the mid.

\section{What we did differently}

The full list is in \texttt{DEVIATIONS.md}. The ones that matter for reading the results:
\begin{itemize}
\item \textbf{Market.} Deribit BTC options instead of end-of-day SPX. These are quoted in
  BTC, which is already a price per unit of forward, so a quote is used as the pure price
  at strike $K/F$. The exchange's own mark prices follow this convention to a median of
  $2\times10^{-6}$ BTC against a tick of $10^{-4}$.
\item \textbf{Filters.} Out-of-the-money options with a bid above the minimum tick, and
  expiries with at least 20 of them. They were fixed from the quotes of the first snapshot.
  There is no volume filter and no distance-to-the-money filter.
\item \textbf{Model strikes.} The paper adds strikes without stating the rule. We add them
  at most $0.1\sqrt{V_j}$ apart, up to $2\sqrt{V_j}$ beyond the outermost quotes.
\item \textbf{Scope.} Model expiries are the quoted expiries. The generalized model, the
  discrete-local-volatility parametrization and extrapolation in time are not implemented.
\end{itemize}

\section{Results}

Synthetic quotes come from a mixture of two Black--Scholes models (3 expiries, 21
strikes). Each Deribit snapshot has 374 quotes on 10 expiries from 2 to 313 days. Errors
are in half-spreads; the claims were written down before the code.

\begin{center}
\begin{tabular}{@{}lp{4.7cm}p{8.3cm}@{}}
\toprule
 & Claim & Result \\
\midrule
E0 & The fit recovers a known surface & Holds for prices. The density is within 1\% of
  the truth at $\eta=0.25$ and within 15\% at $\eta=0.06$, where the quotes do not pin
  it down. \\
E1 & Linear interpolation is too expensive between strikes & Holds. Halfway between
  quoted strikes the linear fit's implied volatility is above the average of its two
  neighbours at 99\% of 364 midpoints; for $\eta=0.25$ at 8--10\%. \\
E2 & $\eta$ trades smoothness against fit & Holds. Density maxima on a short, mid and long
  expiry: 3, 7, 3 at $\eta=0.06$ and 1, 2--3, 1 at $\eta=0.25$, both with every quote
  inside bid/ask. A fit cost appears only at $\eta=0.5$. \\
E3 & A full surface fits within bid/ask in under a second & Holds on a smaller problem:
  100\% of 374 quotes inside bid/ask in 0.3\,s. The paper reports 91.4\% of 1000 SPX
  quotes on 48 expiries. \\
E4 & Extrapolation beyond the quotes is sensible & Not reproduced. Outside the quotes the
  implied volatility depends on $\eta$ and on the snapshot by tens of points. \\
E5 & Strictly arbitrage-free for $\omega=0$ & Holds. No violation above $2\times10^{-11}$
  in price for $\omega=0$; up to $8\times10^{-5}$ in expiry for $\omega=1$. \\
E6 & Mids show arbitrage that bid/ask does not & Holds. The mid-implied density is negative
  at 43 and 46 of 354 strikes, while arbitrage-free prices exist inside bid/ask. \\
\bottomrule
\end{tabular}
\end{center}

\begin{figure}[h]
\centering
\includegraphics[width=\textwidth]{\snap/e2_25DEC26.png}
\caption{The 40-day expiry for three values of $\eta$. Crosses are strikes where three
neighbouring mids imply a negative density.}
\end{figure}

\begin{figure}[h]
\centering
\includegraphics[width=\textwidth]{\snap/e1_between_strikes.png}
\caption{The same expiry near the money. The linear fit bumps upwards between quoted
strikes.}
\end{figure}

\begin{figure}[h]
\centering
\includegraphics[width=\textwidth]{\snap/e4_17NOV26.png}
\caption{The 2-day expiry, two at-the-money standard deviations beyond its quotes. The
three fits agree where there are quotes and nowhere else.}
\end{figure}

\section{What we learned}

\begin{itemize}
\item \textbf{The grid is part of the model.} With model strikes only at the quoted
  strikes, the smooth fit misses synthetic quotes and the linear fit misses the lowest
  one or two real quotes of 7 or 8 of the 10 expiries. Both disappear on a grid finer than the smoothing width
  $\sqrt{\eta V_j}$. The paper leaves this rule unstated.
\item \textbf{The fit is not unique.} The linear program has many optimal solutions once
  every quote is inside bid/ask, because the remaining objective is $10^{-8}$ times the
  distance to the mid. Which one the solver returns moved the synthetic density error at
  $\eta=0.06$ from 7\% to 13\%, and it is what makes the extrapolation unstable.
\item \textbf{``Strictly'' depends on the solver.} At the default tolerance of the solver
  the $\omega=0$ surface violated the conditions by up to $2\times10^{-8}$. At its tightest
  tolerance the violations are rounding error.
\item \textbf{An infeasible program is not an arbitrage.} The hard bid/ask constraint is
  infeasible for the linear model on the quoted strikes and feasible on the denser grid.
\end{itemize}

\section{Limits}

Two snapshots four minutes apart on a Saturday evening are one sample. The claims that
depend on the market (E1--E4, E6) should be rerun on snapshots from different days before
being relied on; capturing one is a single command. Only BTC was used. The paper's problem
size (48 expiries, 1000 quotes) was not reached, so the runtime claim is untested at that
scale. Everything here is regenerated from the frozen snapshots by two scripts in
\texttt{experiments/}.

\end{document}
